\begingroup
\color{resultrevision}
\section{\resultrevision{Results}}
\label{sec:results}
The results below replace the earlier summaries with a verified snapshot of the completed preference-runtime-v2 experiments. Each of Full, no-transfer, and fixed-transfer has 240 group-level training runs and 540 task-specific held-out tests: eight groups, 18 group/scenario targets, and 30 seeds. Across the three methods this gives 720 training runs, 1,620 tests, and 324,000 preference outcomes. Fixed-transfer was the last batch to finish, on 3 November 2026. The analysis reads the completed attempt identified by each completion marker, validates configuration and source hashes, and checks that paired training parameters differ only in the intended transfer fields. Raw training and test files are left unchanged.

For every task/run, testing holds the first training-ranked pair of trees from generation 49 fixed and supplies 200 preferences, with 30 configured simulation replications per preference. It does not choose a different heuristic for each preference or pool several individuals to form a test front. Training and testing use the corrected preference-terminal scale, $100w_1$ and $100w_2$.

\subsection{\resultrevision{Validity and Metric Protocol}}
Full produces two valid objectives for all 108,000 preferences. No-transfer has 192 invalid preferences in five task/run tests, while fixed-transfer has 171 in four. A preference is invalid when either objective contains the inherited $-1$ sentinel. No training seed is removed. Quality measures for the controls use their valid preference subsets, and the failures remain in the validity report. These conditional quality scores do not penalise failed preferences; validity and front quality must therefore be read together.

The inherited evaluator averages each objective over its finite replications. A valid objective mean does not certify that all 30 replications were finite. Thus, the counts concern preference-level output validity, not replication-level failure rates.

The main table uses common scenario-specific min--max bounds and empirical reference fronts pooled over Full MPSLGP, the completed generation-49 PSLGP retests, NSGPII, and MOGP/D. Replacing the PSLGP outcomes changes the bounds in all ten scenarios, so all four methods are recomputed from raw test objectives. The additional two-method comparison constructs a separate reference family from Full and PSLGP only. The ablations retain their existing joint references over Full, no-transfer, and fixed-transfer. The preference-response illustration retains its Full-only bounds. These reference families must not be interchanged. Let $A$ contain one run's valid normalised outcomes, $F=\operatorname{ND}(A)$ its distinct nondominated subset, and $P$ the distinct nondominated union for that scenario under the stated protocol. For Full and PSLGP, $A$ comes from one fixed conditioned heuristic; for NSGPII and MOGP/D, it comes from the saved population Pareto-front heuristics. The indicators are
\begin{equation}
\begin{aligned}
\widetilde{\mathrm{HV}}(A)&=\mathrm{HV}(F;(1.1,1.1))/1.21,\\
\widetilde{\mathrm{IGD}}(A)&=\frac{1}{|P|}\sum_{\mathbf p\in P}\min_{\mathbf a\in A}\|\mathbf p-\mathbf a\|_2,\\
\widetilde{\mathrm{GD}}(A)&=\frac{1}{|A|}\sum_{\mathbf a\in A}\min_{\mathbf p\in P}\|\mathbf a-\mathbf p\|_2.
\end{aligned}
\label{eq:descriptive-test-metrics}
\end{equation}
Both distance indicators use arithmetic means, following the PSLGP definitions~\cite{xu2026pslgp}. The main comparison reports HV and IGD; GD is retained only in the existing transfer-ablation analysis. GD includes every valid preference outcome, not just the nondominated subset. The empirical fronts contain outcomes from the evaluated runs. The analysis archive retains objective bounds, run-level scores, failed-preference records, paired tests, and input hashes.

\subsection{\resultrevision{Main Comparison with Retested PSLGP}}
Table~\ref{tab:latest-main-quality} reports mean(sample standard deviation) over 30 seeds for all 18 group/scenario targets. It replaces the previous PSLGP generation-50 entries with 300 completed generation-49 retests: ten scenarios, 30 seeds and 200 preferences, giving 60,000 preference outcomes. Completion markers, configuration fingerprints, source hashes and row identities are verified. Each retest uses the recorded generation-49 Best Individual from the PSLGP-I training log, with one fixed sequencing/routing pair over all preferences and 30 simulation replications per preference. Two preference vectors are invalid and excluded conditionally; no seed is removed. Independent baseline values repeat where a scenario occurs in several Full groups; these are not additional experiments.

Under the new common references, Full has less favourable mean HV and IGD than PSLGP and NSGPII in all 18 targets. Against MOGP/D, Full has higher mean HV in 3 targets and lower mean IGD in 11. For scenario 1 in R1, Full gives HV 0.938 (0.012) and IGD 0.032 (0.009), while PSLGP gives 0.941 (0.013) and 0.031 (0.011). For scenario 4 in R2, Full instead gives HV 0.700 (0.060) and IGD 0.190 (0.057), against PSLGP at 0.750 (0.047) and 0.152 (0.039). The results do not establish a general quality advantage for Full.

Changing the PSLGP source changes the common reference data, so changes in the other columns are not new outcomes from their heuristics. A diagnostic calculation holds the previous four-method bounds and fronts fixed: the new PSLGP generation-49 means have lower HV and higher IGD than the old generation-50 means in all ten scenarios. However, the old analysis used generation-50 Top-N rank 0, whereas the new tests use the recorded generation-49 Best Individual because no generation-49 Top-N block was saved. Both generation and selection policy have changed; the observed difference cannot be attributed solely to one fewer generation.

Full and PSLGP now use actual generation 49, but NSGPII and MOGP/D retain their verified generation-50 population Pareto fronts, not external archives or a union of earlier fronts. Their raw outcomes are unchanged. Full shares 1,000 individuals across a group while each independent baseline task uses 1,000. The 50\% or 66.7\% reduction in configured group-level population slots is not a measured speedup or equal-accuracy guarantee. This remains a descriptive comparison, not a fully generation-, selection- or population-budget-matched experiment; no cross-method significance marks are claimed.

\begin{table*}[t]
\centering
\color{resultrevision}
\caption{\resultrevision{Four-method HV and IGD comparison with retested PSLGP generation 49: mean(sample standard deviation) over 30 seeds per target, shown to three decimal places. All four methods use common scenario-specific bounds and empirical reference fronts rebuilt from their raw outcomes; published indicator summaries are not copied. MPSLGP and PSLGP use actual generation 49, while NSGPII and MOGP/D retain generation-50 population Pareto fronts. PSLGP denotes PSLGP-I and uses the recorded generation-49 Best Individual, not the previous generation-50 Top-N rank 0. Two invalid PSLGP preference outcomes are excluded conditionally; all seeds remain. Higher HV and lower IGD are better; bold marks the best unrounded mean, not statistical significance.}}
\label{tab:latest-main-quality}
\footnotesize
\setlength{\tabcolsep}{1.5pt}
\newcommand{\PSLGPFortyNineRows}{%
R1 & 1 & 0.938(0.012) & 0.941(0.013) & \textbf{0.976(0.007)} & 0.941(0.009) & 0.032(0.009) & 0.031(0.011) & \textbf{0.011(0.003)} & 0.033(0.005) \\
R1 & 2 & 0.935(0.008) & 0.941(0.012) & \textbf{0.980(0.007)} & 0.941(0.008) & 0.037(0.006) & 0.034(0.009) & \textbf{0.014(0.004)} & 0.039(0.008) \\
R2 & 3 & 0.946(0.011) & 0.948(0.009) & \textbf{0.982(0.006)} & 0.953(0.007) & 0.029(0.008) & 0.027(0.007) & \textbf{0.011(0.003)} & 0.030(0.006) \\
R2 & 4 & 0.700(0.060) & 0.750(0.047) & \textbf{0.923(0.025)} & 0.780(0.029) & 0.190(0.057) & 0.152(0.039) & \textbf{0.052(0.014)} & 0.156(0.025) \\
R3 & 6 & 0.726(0.051) & 0.756(0.037) & \textbf{0.913(0.035)} & 0.726(0.034) & 0.200(0.048) & 0.169(0.033) & \textbf{0.067(0.020)} & 0.209(0.028) \\
R3 & 7 & 0.912(0.032) & 0.922(0.013) & \textbf{0.976(0.010)} & 0.912(0.013) & 0.059(0.035) & 0.050(0.012) & \textbf{0.015(0.005)} & 0.059(0.010) \\
R4 & 1 & 0.938(0.012) & 0.941(0.013) & \textbf{0.976(0.007)} & 0.941(0.009) & 0.033(0.008) & 0.031(0.011) & \textbf{0.011(0.003)} & 0.033(0.005) \\
R4 & 3 & 0.942(0.011) & 0.948(0.009) & \textbf{0.982(0.006)} & 0.953(0.007) & 0.031(0.008) & 0.027(0.007) & \textbf{0.011(0.003)} & 0.030(0.006) \\
R4 & 5 & 0.899(0.016) & 0.912(0.012) & \textbf{0.979(0.010)} & 0.930(0.013) & 0.063(0.013) & 0.052(0.008) & \textbf{0.016(0.005)} & 0.056(0.008) \\
R5 & 8 & 0.677(0.067) & 0.708(0.050) & \textbf{0.908(0.033)} & 0.676(0.050) & 0.220(0.059) & 0.196(0.039) & \textbf{0.075(0.016)} & 0.252(0.036) \\
R5 & 9 & 0.886(0.015) & 0.899(0.021) & \textbf{0.973(0.010)} & 0.912(0.012) & 0.068(0.012) & 0.060(0.020) & \textbf{0.018(0.008)} & 0.071(0.010) \\
H1 & 1 & 0.940(0.013) & 0.941(0.013) & \textbf{0.976(0.007)} & 0.941(0.009) & 0.034(0.016) & 0.031(0.011) & \textbf{0.011(0.003)} & 0.033(0.005) \\
H1 & 10 & 0.661(0.072) & 0.673(0.070) & \textbf{0.923(0.028)} & 0.702(0.034) & 0.216(0.068) & 0.210(0.062) & \textbf{0.059(0.021)} & 0.240(0.034) \\
H2 & 5 & 0.903(0.014) & 0.912(0.012) & \textbf{0.979(0.010)} & 0.930(0.013) & 0.059(0.011) & 0.052(0.008) & \textbf{0.016(0.005)} & 0.056(0.008) \\
H2 & 6 & 0.733(0.040) & 0.756(0.037) & \textbf{0.913(0.035)} & 0.726(0.034) & 0.193(0.036) & 0.169(0.033) & \textbf{0.067(0.020)} & 0.209(0.028) \\
H3 & 1 & 0.940(0.010) & 0.941(0.013) & \textbf{0.976(0.007)} & 0.941(0.009) & 0.031(0.006) & 0.031(0.011) & \textbf{0.011(0.003)} & 0.033(0.005) \\
H3 & 5 & 0.895(0.016) & 0.912(0.012) & \textbf{0.979(0.010)} & 0.930(0.013) & 0.065(0.014) & 0.052(0.008) & \textbf{0.016(0.005)} & 0.056(0.008) \\
H3 & 8 & 0.657(0.060) & 0.708(0.050) & \textbf{0.908(0.033)} & 0.676(0.050) & 0.235(0.056) & 0.196(0.039) & \textbf{0.075(0.016)} & 0.252(0.036) \\
}

\begin{tabular}{@{}lc*{8}{c}@{}}
\toprule
 & & \multicolumn{4}{c}{HV $\uparrow$} & \multicolumn{4}{c}{IGD $\downarrow$} \\
\cmidrule(lr){3-6}\cmidrule(lr){7-10}
Group & Scenario & Full MPSLGP & PSLGP & NSGPII & MOGP/D & Full MPSLGP & PSLGP & NSGPII & MOGP/D \\
\midrule
\PSLGPFortyNineRows
\bottomrule
\end{tabular}
\end{table*}

\subsection{\resultrevision{Direct Comparison with PSLGP}}
Table~\ref{tab:latest-two-method-quality} compares only Full MPSLGP and the retested PSLGP-I heuristics. The calculation uses the same 540 Full task/run tests and 300 PSLGP tests at actual generation 49, but rebuilds each scenario's bounds and empirical reference front from only these two methods. No NSGPII or MOGP/D outcomes enter this reference pool. All 30 seeds per target are retained; PSLGP's two invalid preferences are excluded conditionally. Table~\ref{tab:latest-main-quality} remains the four-method reference comparison.

Removing the other methods changes the normalisation bounds in all ten scenarios, so this is not obtained by deleting columns from Table~\ref{tab:latest-main-quality}. Under the two-method references, Full still has lower mean HV and higher mean IGD than PSLGP in all 18 targets. For scenario 1 in R1, Full gives HV 0.967 (0.013) and IGD 0.019 (0.008), compared with PSLGP at 0.970 (0.014) and 0.017 (0.011). For scenario 4 in R2, the corresponding values are 0.752 (0.062) and 0.160 (0.060) for Full, against 0.822 (0.057) and 0.108 (0.050) for PSLGP. Differences from the four-method table reflect the reference scope, not new scheduling outcomes. No statistical superiority or equivalence is inferred from these mean comparisons.

Both methods now use generation-49 outputs and one conditioned heuristic per task/run. However, Full selects task-specific Top-N rank 0 while PSLGP uses its recorded Best Individual, and Full shares 1,000 individuals across tasks while PSLGP trains each task with 1,000. The restricted comparison removes the other algorithms from metric construction but does not remove these selection and population-budget differences. It therefore does not establish an accuracy advantage for MPSLGP.

\begin{table*}[t]
\centering
\color{resultrevision}
\caption{\resultrevision{Two-method HV and IGD comparison: mean(sample standard deviation) over 30 seeds per group/scenario, shown to three decimal places. Only Full MPSLGP and retested PSLGP-I contribute to the scenario-specific normalisation bounds and empirical reference fronts; NSGPII and MOGP/D outcomes are excluded from this calculation. Both methods use actual generation 49 and one fixed dual-tree heuristic across 200 preferences. MPSLGP selects task-specific Top-N rank 0; PSLGP uses its recorded Best Individual. Higher HV and lower IGD are better; bold marks the best unrounded mean, not statistical significance. Two invalid PSLGP preferences are excluded conditionally; all seeds remain. This is not a population-budget- or selection-matched experiment.}}
\label{tab:latest-two-method-quality}
\footnotesize
\setlength{\tabcolsep}{6pt}
\newcommand{\MPSLGPPSLGPOnlyRows}{%
R1 & 1 & 0.967(0.013) & \textbf{0.970(0.014)} & 0.019(0.008) & \textbf{0.017(0.011)} \\
R1 & 2 & 0.959(0.008) & \textbf{0.964(0.012)} & 0.024(0.006) & \textbf{0.019(0.007)} \\
R2 & 3 & 0.966(0.012) & \textbf{0.969(0.009)} & 0.019(0.008) & \textbf{0.017(0.007)} \\
R2 & 4 & 0.752(0.062) & \textbf{0.822(0.057)} & 0.160(0.060) & \textbf{0.108(0.050)} \\
R3 & 6 & 0.818(0.058) & \textbf{0.853(0.043)} & 0.130(0.037) & \textbf{0.103(0.025)} \\
R3 & 7 & 0.940(0.033) & \textbf{0.951(0.013)} & 0.033(0.033) & \textbf{0.024(0.008)} \\
R4 & 1 & 0.967(0.012) & \textbf{0.970(0.014)} & 0.019(0.008) & \textbf{0.017(0.011)} \\
R4 & 3 & 0.962(0.011) & \textbf{0.969(0.009)} & 0.020(0.008) & \textbf{0.017(0.007)} \\
R4 & 5 & 0.935(0.017) & \textbf{0.949(0.013)} & 0.037(0.014) & \textbf{0.028(0.008)} \\
R5 & 8 & 0.736(0.091) & \textbf{0.768(0.080)} & 0.164(0.080) & \textbf{0.141(0.061)} \\
R5 & 9 & 0.929(0.016) & \textbf{0.943(0.023)} & 0.040(0.011) & \textbf{0.032(0.020)} \\
H1 & 1 & 0.969(0.013) & \textbf{0.970(0.014)} & 0.021(0.016) & \textbf{0.017(0.011)} \\
H1 & 10 & 0.659(0.079) & \textbf{0.695(0.082)} & 0.175(0.073) & \textbf{0.150(0.078)} \\
H2 & 5 & 0.939(0.014) & \textbf{0.949(0.013)} & 0.034(0.010) & \textbf{0.028(0.008)} \\
H2 & 6 & 0.826(0.045) & \textbf{0.853(0.043)} & 0.122(0.037) & \textbf{0.103(0.025)} \\
H3 & 1 & 0.969(0.010) & \textbf{0.970(0.014)} & 0.017(0.005) & \textbf{0.017(0.011)} \\
H3 & 5 & 0.930(0.017) & \textbf{0.949(0.013)} & 0.040(0.015) & \textbf{0.028(0.008)} \\
H3 & 8 & 0.709(0.080) & \textbf{0.768(0.080)} & 0.185(0.067) & \textbf{0.141(0.061)} \\
}

\begin{tabular}{@{}lc*{4}{c}@{}}
\toprule
 & & \multicolumn{2}{c}{HV $\uparrow$} & \multicolumn{2}{c}{IGD $\downarrow$} \\
\cmidrule(lr){3-4}\cmidrule(lr){5-6}
Group & Scenario & Full MPSLGP & PSLGP & Full MPSLGP & PSLGP \\
\midrule
\MPSLGPPSLGPOnlyRows
\bottomrule
\end{tabular}
\end{table*}


\subsection{\resultrevision{Transfer Ablations with Common References}}
Figures~\ref{fig:latest-ablation-quality} and~\ref{fig:latest-ablation-effects} compare Full, no-transfer, and fixed-transfer after recomputing all their indicators with joint scenario-specific references. Against no-transfer, Full has better mean HV, IGD, and GD in 15, 14, and 17 of the 18 targets. Against fixed-transfer, the corresponding counts are 7, 7, and 9. Transfer presence has a broadly favourable mean pattern relative to local mating, but the present adaptive controller is not consistently better than fixed transfer.

For scenario 3 in R4, Full gives joint HV 0.9689 (0.0092), IGD 0.0195 (0.0062), and GD 0.0143 (0.0045). No-transfer gives 0.9621 (0.0093), 0.0233 (0.0079), and 0.0183 (0.0050). All three means favour Full in this target. For scenario 1 in R1, fixed-transfer instead has joint HV 0.9439 (0.0153), compared with Full at 0.9318 (0.0229), and also has lower mean IGD and GD. Both examples remain in the analysis.

None of the 108 exact paired signed-rank tests survives Holm correction at 0.05. The smallest unadjusted $p$ is 0.000730 for Full versus no-transfer GD on R4/scenario~3; its adjusted value is 0.078797. A non-significant result is not evidence of equivalence. The mean patterns support further study of transfer but do not establish adaptive-transfer superiority. These tests use 30 paired training seeds for each fixed target; the 200 preferences and 30 replications are not additional independent evolutionary runs.

The map displays relative changes in target means. For an indicator mean $q$ and control $c$, the signed percentage is
\begin{equation}
\Delta_q=100\,\sigma_q\frac{\bar q_{\mathrm{Full}}-\bar q_c}{\bar q_c},
\qquad
\sigma_q=\begin{cases}1,&q=\mathrm{HV},\\-1,&q=\mathrm{IGD},\mathrm{GD}.\end{cases}
\end{equation}
Large percentage gains in GD can reflect a few poor control runs as well as a change in typical performance. Means, dispersion, and paired tests should therefore be considered together.

\begin{figure*}[t]
\centering
\includegraphics[width=\textwidth]{results/latest-20261103/ablation-joint-quality.pdf}
\caption{\resultrevision{Completed transfer ablations under common reference data. Points and bars are means and sample standard deviations over 30 paired seeds. Full, no-transfer, and fixed-transfer share scenario-specific bounds and empirical fronts. Control quality scores are conditional on valid preferences; no seed is removed.}}
\label{fig:latest-ablation-quality}
\end{figure*}

\begin{figure*}[t]
\centering
\includegraphics[width=\textwidth]{results/latest-20261103/ablation-effect-map.pdf}
\caption{\resultrevision{Relative changes in target means for Full against each ablation. Positive values favour Full; negative values favour the control. The map describes effect direction and size, not significant wins. No test survives the 108-comparison Holm family.}}
\label{fig:latest-ablation-effects}
\end{figure*}

\subsection{\resultrevision{Training-Phase Transfer Behaviour}}
Figure~\ref{fig:latest-transfer-training} shows the actual Full probability tables across 30 runs per group. The curves average total cross-task probabilities over receiving tasks and preference regions; they are not observed mating frequencies. The two-task groups start near 0.380 and end at 0.367--0.368, while R4 and H3 move from 0.258 to about 0.252. The controller changes its stored probabilities, but the aggregate adjustment is modest under these settings.

Fixed-transfer uses probabilities 0.8 for two tasks and 0.4 for three tasks, whereas the no-transfer probability is zero. Full's no-transfer reference therefore changes the realised transfer frequency even when the configured budget is shared. The adaptive-versus-fixed ablation combines this frequency effect with utility-based donor selection. It cannot isolate improved donor selection, the value of preference regions, or a causal reduction in negative transfer. A frequency-matched fixed controller and a $K=1$ controller remain needed to separate these effects.

\begin{figure*}[t]
\centering
\includegraphics[width=\textwidth]{results/latest-20261103/training-transfer-probabilities.pdf}
\caption{\resultrevision{Training-phase transfer probabilities from all 240 latest Full runs. Each curve is the mean of 30 run-level task/region averages after utility updates; shading is sample standard deviation. Dashed lines are the configured fixed-transfer mating rates and dotted lines denote no transfer. These stored probabilities are not empirical mating frequencies.}}
\label{fig:latest-transfer-training}
\end{figure*}

\subsection{\resultrevision{Validity and Preference Response}}
Figure~\ref{fig:latest-validity-response} separates output validity from response to preferences. Full has no invalid preference vectors, compared with 0.178\% for no-transfer and 0.158\% for fixed-transfer. These failures are concentrated in five and four task/run tests, respectively. The percentages should not be treated as thousands of independent failure observations because preferences from one trained heuristic are dependent.

Full has 48 single-outcome heuristics among 540 (8.89\%); no-transfer has 39 (7.22\%) and fixed-transfer 57 (10.56\%). Full is therefore more reliable at producing valid outputs in this snapshot, but it does not have the lowest observed single-outcome count. Mean distinct outcome counts are 170.53, 170.10, and 166.21, respectively, whereas mean within-heuristic nondominated counts are only 5.65, 5.69, and 5.37. A preference-conditioned representation can produce many responses without producing equally many useful trade-offs.

Figure~\ref{fig:latest-preference-response} fixes one common illustrative seed, 11, for the R4 and H3 task heuristics. It is selected as the lower median of the equal-group mean within-task IGD ranks over seeds 0--29, rather than the best-performing run. The curves show all 200 preferences. Several responses are irregular; both objectives can move in the same direction over part of the range. The illustration does not establish monotonic preference alignment or front coverage for all runs.

\begin{figure*}[t]
\centering
\includegraphics[width=\textwidth]{results/latest-20261103/validity-and-response.pdf}
\caption{\resultrevision{Preference-level invalid counts, single-outcome heuristic counts, and mean distinct nondominated outcomes for the completed methods. Invalid percentages use 108,000 preferences per method; single-outcome percentages use 540 trained task heuristics. Validity and responsiveness are different properties.}}
\label{fig:latest-validity-response}
\end{figure*}

\begin{figure*}[t]
\centering
\includegraphics[width=\textwidth]{results/latest-20261103/latest-preference-response.pdf}
\caption{\resultrevision{Latest Full preference responses for the co-trained tasks in R4 and H3, seed 11. Each panel uses one fixed pair of trees and all 200 preferences. Values are normalised by the Full-only scenario bounds, not the three-method ablation bounds. The seed is selected by the stated median-rank rule.}}
\label{fig:latest-preference-response}
\end{figure*}

\subsection{\resultrevision{Raw-Objective Attainment and Scale Sensitivity}}
Changing the pooled reference range changes the absolute indicators even when the Full test outcomes remain fixed. R3/scenario~7, for example, has Full-only HV 0.731 but joint HV 0.984. The joint normalisation includes finite control outcomes up to $WF_{\max}=29044.5$ and $T_{\max}=10600.8$, which compress the region occupied by most Full outcomes. This difference is not a new improvement in the learned heuristic. Full-only and joint-reference scores must not be interchanged, and isolated near-one HV values do not establish strong Pareto-front approximation.

Figure~\ref{fig:latest-attainment} therefore also uses raw objective space. For each Full group/task/seed, one fixed heuristic supplies its nondominated preference outcomes. The empirical attainment at a target point is the fraction of the 30 seed fronts that weakly dominate that point. Median boundaries and the 25--75\% envelopes compare R4 and H3 on their shared scenarios 1 and 5. No outcomes are removed as plotting outliers. These are comparisons of training groups within Full, not raw-front comparisons against PSLGP. A lower-left boundary is favourable, and crossings mean that no grouping is uniformly preferable across trade-offs.

\begin{figure*}[t]
\centering
\includegraphics[width=\textwidth]{results/latest-20261103/latest-empirical-attainment.pdf}
\caption{\resultrevision{Raw-objective empirical attainment for latest Full runs in R4 and H3 on shared scenarios 1 and 5. Solid staircases show 50\% attainment, and shading spans 25--75\%, requiring attainment by at least 8, 15, and 23 of 30 run fronts. Curves are exact right-continuous boundaries; no pooled normalisation or outlier removal is used.}}
\label{fig:latest-attainment}
\end{figure*}

\subsection{\resultrevision{Training Cost and Heuristic Size}}
Figure~\ref{fig:latest-cost} records whole-group training time once per method/group/seed, rather than summing the same run time over its task rows. Full's mean group time ranges from 80.42 to 146.02 minutes; no-transfer ranges from 78.36 to 105.96 and fixed-transfer from 82.56 to 106.29. Full R1 has a large standard deviation of 118.07 minutes. Batch concurrency and run conditions differ, so these observations cannot establish a hardware-independent speedup over either ablation or the published PSLGP experiment.

The mean total size of the tested dual-tree heuristics is 82.82 nodes for Full, 82.72 for no-transfer, and 82.43 for fixed-transfer. Rule sizes are thus similar in aggregate. Neither lower run time nor a small tree is sufficient evidence of preference responsiveness or approximation quality. The present results support a distinction between completion reliability, preference response, and front quality; the adaptive method does not dominate all three criteria over both completed controls.

\begin{figure*}[t]
\centering
\includegraphics[width=\textwidth]{results/latest-20261103/runtime-and-complexity.pdf}
\caption{\resultrevision{Latest training costs and tested heuristic sizes. Training time is whole-group wall time, with mean and sample standard deviation over 30 runs. Heuristic size is the mean sum of sequencing and routing nodes over all tested task heuristics in the group. Concurrent batch execution limits causal interpretation of timing differences.}}
\label{fig:latest-cost}
\end{figure*}
\endgroup